% Supporting analysis of the public implementations. No main-paper proof is
% moved here, and this file has no dependency on the private manuscript.
\newtheorem{auditproposition}{Proposition}
\renewcommand{\theauditproposition}{S.A\arabic{auditproposition}}

\section{Learning-rule audit and additional supporting proofs}
\label{sec:learning_rule_audit}

This section audits the retained coefficient-based contact, Burgers, and
bounded Fourier studies; the current direct-field main experiments are
documented separately in Section~\ref{sec:direct_main}.
They estimate coefficients of prescribed implicit relations
and obtain their evaluated context partials by differentiation. Their adaptive
updates receive observed features, not target Jacobian entries. This is
coefficient estimation with a structured differential evaluation. It is
not the direct all-entry learning mechanism specified in
Section~\ref{sec:scope}: physical
coefficients, feature families, residual levels, and selected sensing channels
are supplied. These studies implement constrained sampled
coefficient updates; they do not numerically integrate the unrestricted
row-Stiefel learning flow. The distinctions below specify what the numerical
evidence can support. The retained normalized Fourier experiment is audited
separately in Section~\ref{sec:normalized_fourier_results}; it adapts all
seven coefficients and numerically integrates the coupled normalized law,
but still derives its partials from a supplied residual family.

\subsection{Exact update, scaling, and current-candidate memory}
Write the physical free coefficients as $c=D_s\theta$, where $D_s$ is the
positive diagonal matrix of optimization scales. For observed free features
$a_{n,\ell}$ and a supplied contribution $b_{n,\ell}$, the implemented
implicit feature row is
\begin{equation}
 v_{n,\ell}=\col(D_sa_{n,\ell},b_{n,\ell}),\qquad
 q_\theta=\col(\theta,1).
 \label{eq:audit_channels}
\end{equation}
Thus $q_\theta^{\mathsf T}v_{n,\ell}=c^{\mathsf T}a_{n,\ell}+b_{n,\ell}$;
scaling changes optimization coordinates, not the physical residual.
Let $N_n$ be the number of simultaneous spatial samples and let
$\beta=\exp(-h/\tau_m)$. After scoring the new sample with $\theta_{n-1}$,
the code sets
\begin{align}
 R_n&=\beta R_{n-1}+\frac{1-\beta}{N_n}
       \sum_{\ell=1}^{N_n}v_{n,\ell}v_{n,\ell}^{\mathsf T},\\
 H_n&=(R_n)_{1:p,1:p},\qquad g_n=(R_n)_{1:p,p+1}.
\end{align}
Starting from $R_0=0$ and $\theta_n^{(0)}=\theta_{n-1}$, it sets
\begin{align}
 \alpha_n&=\frac{1}{\max(\tr H_n,10^{-12})},\\
 \theta_n^{(j+1)}&=\Pi_{\mathcal W}
 \bigl(\theta_n^{(j)}-\alpha_n(H_n\theta_n^{(j)}+g_n)\bigr).
 \label{eq:audit_exact_update}
\end{align}
The output is $\theta_n=\theta_n^{(M)}$ and $\mathcal W$ is the scaled
coefficient box. Contact uses $M=48$;
Bounded Fourier and Burgers use $M=8$. Burgers stores its unscaled moment with
the fixed channel first; permuting that channel and applying $D_s$ gives
exactly the same recurrence. Material and secant observations begin only
when the specified observation window is full. Contact's kinematic rows
start accumulating at the first sample. At the terminal sample the memory
may still be updated, but the coefficients and physical state are not.

\begin{auditproposition}[One candidate over the entire stored history]
\label{prop:audit_memory}
For any sequence of observed feature rows and any candidate $\vartheta$,
including a candidate never used when the data were collected,
\begin{equation}
 \frac12q_\vartheta^{\mathsf T}R_nq_\vartheta
 =\frac{1-\beta}{2}\sum_{i=1}^{n}\frac{\beta^{n-i}}{N_i}
   \sum_{\ell=1}^{N_i}(q_\vartheta^{\mathsf T}v_{i,\ell})^2.
 \label{eq:audit_memory_identity}
\end{equation}
Its free gradient is $H_n\vartheta+g_n$. Each inner step in
(\ref{eq:audit_exact_update}) decreases this fixed objective by at least
$\|\theta_n^{(j+1)}-\theta_n^{(j)}\|^2/(2\alpha_n)$.
\end{auditproposition}
\begin{IEEEproof}
Unrolling the linear recursion gives
$R_n=(1-\beta)\sum_{i=1}^n\beta^{n-i}N_i^{-1}
\sum_\ell v_{i,\ell}v_{i,\ell}^{\mathsf T}$.
Multiplication by the same $q_\vartheta$ on both sides proves
(\ref{eq:audit_memory_identity}); differentiation gives the stated
gradient. Let $d=\theta^+-\theta$ for one projected step. Projection onto
the convex box gives $(H_n\theta+g_n)^{\mathsf T}d\le-\|d\|^2/\alpha_n$.
The exact quadratic expansion then gives
\begin{align*}
 \mathcal J_n(\theta+d)-\mathcal J_n(\theta)
 &\le-\|d\|^2/\alpha_n+d^{\mathsf T}H_nd/2\\
 &\le-\|d\|^2/(2\alpha_n),
\end{align*}
because $H_n\succeq0$ and
$\alpha_n\lambda_{\max}(H_n)\le\alpha_n\tr H_n\le1$.
\end{IEEEproof}
The identity holds for dependent, overlapping, closed-loop observations;
independence is unnecessary for this algebra. It proves that old residuals
are reevaluated at the current candidate through sufficient statistics.
It does not prove statistical consistency, physical excitation, or monotone
energy when a new observation changes $R_n$.

\subsection{What is observed and what is learned}
\subsubsection{Contact}
The material relation is
\begin{equation}
 f_c=\dot\lambda-s^{\mathsf T}\beta(x)+ay+\rho\lambda
       -r^{\mathsf T}\beta'(x)v_x+bv_y=0,
 \label{eq:audit_contact_relation}
\end{equation}
with learned $c=\col(s,a,r,b)$. Its partials, in the declared context
$(x,y,\lambda,v_x,v_y,\dot\lambda,u_1,u_2)$, are
\begin{equation}
 F_c=\bigl[-s^{\mathsf T}\beta'(x)-r^{\mathsf T}\beta''(x)v_x,
 a,\rho,-r^{\mathsf T}\beta'(x),b,1,0,0\bigr].
 \label{eq:audit_contact_partials}
\end{equation}
The two additional rows come from $v-Bu=0$; their velocity identity block
is supplied and their input block $-B$ is learned. Both material vectors
$s,r$, both material scalars $a,b$, and all four entries of $B$ adapt.
The constant $\rho=4$, the unit coefficient of $\dot\lambda$, the unit
velocity block, and the functional form of $\beta$ are supplied.

Let $\Delta$ be the material observation window and let $\delta$ denote
its endpoint difference. Integrating (\ref{eq:audit_contact_relation})
with one fixed candidate gives the observed material incompatibility
\begin{equation}
 \begin{split}
 \xi_c=\Delta^{-1}\biggl[&\delta\lambda-s^{\mathsf T}\!\int\beta(x)\dd t
 +a\!\int y\dd t\\
 &+\rho\!\int\lambda\dd t-r^{\mathsf T}\delta\beta(x)+b\delta y\biggr].
 \end{split}
 \label{eq:audit_contact_integral}
\end{equation}
The implementation uses trapezoidal quadrature and measured force for
$\lambda$; it does not train on clean $\dot\lambda$ or on a target $F_c$.
The kinematic learner uses the instantaneous measured residual $v_m-Bu$.
Consequently these two contact updates use calibrated integral-level and
velocity-level relations, respectively. They are not measurements of the
raw-context differential incompatibility $F\dot z$. At the material
change, a window can contain two different physical coefficient sets;
such windows are retained and marked, rather than assumed exactly
compatible with either fixed coefficient vector.

The controller estimates the surface slope from the noisy two-component
force measurement:
\begin{equation*}
 \widehat h'_m=-f_{m,x}/\max(f_{m,y},0.2).
\end{equation*}
The simulator generates
that force through the true surface, as a physical sensor would respond to
the environment. The controller does not receive the true geometric
coefficient vector. Its curvature estimate uses the learned ratio $s/a$.
Thus geometry is partly observed through force direction and partly
inferred through the learned relation. The true surface also fixes the
initial physical contact and the evaluator's indentation diagnostic.
These are separate from the adaptive update.

\subsubsection{Bounded Fourier comparison}
For $h(\phi)=(1,\cos\phi,\sin\phi)^{\mathsf T}$,
the learned six coefficients define
\begin{align}
 \rho_c(\phi)&=c_\rho^{\mathsf T}h(\phi),&
 \kappa_c(\phi)&=c_\kappa^{\mathsf T}h(\phi),\\
 f_c(x,v,u,\phi)&=\rho_c(\phi)v+\kappa_c(\phi)x-u.
\end{align}
The evaluated derivative is
\begin{equation}
 F_c=\bigl[\kappa_c,\rho_c,-1,
 c_\rho^{\mathsf T}h'(\phi)v+c_\kappa^{\mathsf T}h'(\phi)x\bigr].
 \label{eq:audit_fourier_partials}
\end{equation}
The state, rate, and phase columns therefore depend on learned coefficients;
only the input partial $-1$ is fixed. The rate measurement is noisy; the
state and phase are exact in this model. The phase law $\phi=4t$ and the
harmonic basis are supplied, so the experiment does not learn a frequency
or discover the basis. The displayed phase partial is with respect to
$\phi$, not $t$; the clock contribution along the trajectory is multiplied
by $\dot\phi=4$.

The learner differences the complete observed feature vector
$(h(\phi)v_m,h(\phi)x,-u)$ over the declared window, with the same current
coefficient vector at both endpoints. Thus the phase-state and phase-rate
product terms are retained. A $K=0$ shadow receives the identical stream
but is restricted to two constant coefficients; it has the same memory
mechanism. It is a basis ablation, not an independent estimator of matched
capacity. The separate $K=0$ controller generates its own trajectory.

\subsubsection{Burgers}
At each grid point, with context $(w,w_s,w_{ss},u,w_t)$,
\begin{align}
 f_c&=w_t+c_a w w_s-\nu w_{ss}-bu,\\
 F_c&=(c_a w_s,c_a w,-\nu,-b,1).
 \label{eq:audit_burgers_partials}
\end{align}
The three coefficients $c_a,\nu,b$ adapt. Four of the five evaluated
partials depend on those coefficients; the $w_t$ partial is supplied.
The learner differences complete features $(w_{t,m},ww_s,-w_{ss},-u)$,
including the nonlinear product as a whole. State values and their
spectral spatial derivatives are exact sensing channels; only $w_t$ is
noisy. The plant and learner use the same collocation product and spatial
operators, with full-field sensing and actuation. This is a matched
semidiscrete-model experiment, not independent continuous-PDE validation
or evidence for sparse sensors, unknown spatial operators, or model mismatch.

\subsection{Complete-feature secants and the derivative being estimated}
\begin{auditproposition}[Fixed-candidate secant identity]
\label{prop:audit_secant}
Let $f_\theta(z)=Q_\theta\zeta(z)$ be continuously differentiable on the
line segment between $z_-$ and $z_+$ and fix $\theta$. For $\Delta>0$,
\begin{align}
 Q_\theta\frac{\zeta(z_+)-\zeta(z_-)}{\Delta}
 &=\frac{f_\theta(z_+)-f_\theta(z_-)}{\Delta}\\
 &=\frac{1}{\Delta}\int_0^1 F_\theta(z_-+s\delta z)\delta z\dd s,
 \label{eq:audit_secant_identity}
\end{align}
where $\delta z=z_+-z_-$. If $F_\theta$ is $K$-Lipschitz on that segment,
the difference between this secant and
$F_\theta(z_-)\delta z/\Delta$ is at most
$K\|\delta z\|^2/(2\Delta)$.
\end{auditproposition}
\begin{IEEEproof}
Linearity gives the first equality. Apply the fundamental theorem of
calculus to $s\mapsto f_\theta(z_-+s\delta z)$ for the second. Subtract
$F_\theta(z_-)\delta z/\Delta$ inside the integral and use
$\|F_\theta(z_-+s\delta z)-F_\theta(z_-)\|\le Ks\|\delta z\|$;
integration of $s$ over $[0,1]$ gives the bound.
\end{IEEEproof}
This result includes every complete-product and clock term without
assuming that held commands are temporally smooth. The segment is an
algebraic interpolation of context endpoints, not a claim that the plant
followed that segment. With noisy endpoints or command changes, the
observed secant is not an exact instantaneous sample of $F\dot z$.

In all three models $F_\theta=D_zf_\theta$ holds with the candidate fixed.
If both the state and coefficients vary, the different identity
\begin{equation}
 \frac{\dd}{\dd t}f_{\theta(t)}(z(t))
 =F_{\theta(t)}(z(t))\dot z(t)
   +\partial_\theta f_{\theta(t)}(z(t))\dot\theta(t)
 \label{eq:audit_adapting_residual}
\end{equation}
contains the parameter term. The code does not difference residual values
evaluated with old coefficients at one endpoint and new coefficients at
the other, so it does not inadvertently insert or omit this adaptation
term in a fixed-candidate derivative. Compatibility alone also cannot
identify an additive residual level: $f$ and $f+C$ have the same
differential and complete-feature differences. The direct readouts use
the explicitly supplied zero level.

\subsection{A conditional tracking bound for the actual sampled update}
The following result applies to the implemented box-constrained quadratic
update. It makes no identification claim when excitation is absent and
does not import a Stiefel-flow convergence rate.

\begin{auditproposition}[Optimization tracking and physical bias]
\label{prop:audit_tracking}
At each observation $n$, suppose $\mathcal W$ is a fixed nonempty closed
convex set, $\mu_n I\preceq H_n\preceq L_nI$ with $\mu_n>0$, and
$0<\alpha_n\le1/L_n$. Let $\theta_n^\circ$ be the unique minimizer
of $\mathcal J_n$ over $\mathcal W$, and let $M_n\ge1$ projected steps begin
at $\theta_{n-1}$. Define
\begin{equation}
 r_n=(1-\alpha_n\mu_n)^{M_n},\quad
 d_n=\|\theta_n^\circ-\theta_{n-1}^\circ\|,\quad
 e_n=\|\theta_n-\theta_n^\circ\|.
\end{equation}
Then $e_n\le r_n(e_{n-1}+d_n)$. If $r_n\le\bar r<1$ and
$d_n\le\bar d$ on an interval of $n$ observations,
\begin{equation}
 e_n\le\bar r^{\,n}e_0+
       \frac{\bar r(1-\bar r^{\,n})}{1-\bar r}\bar d.
 \label{eq:audit_optimizer_tracking}
\end{equation}
For any feasible physical coefficient target $\theta_n^*$,
\begin{equation}
 \|\theta_n^\circ-\theta_n^*\|
 \le\frac{\|\nabla\mathcal J_n(\theta_n^*)\|}{\mu_n}.
 \label{eq:audit_physical_bias}
\end{equation}
The indices can be restarted after a warm-up interval on which positive
curvature was not available.
\end{auditproposition}
\begin{IEEEproof}
The minimizer satisfies the projection fixed-point identity
$\theta_n^\circ=\Pi_{\mathcal W}(\theta_n^\circ-
\alpha_n\nabla\mathcal J_n(\theta_n^\circ))$.
Projection is nonexpansive, while the difference of the two unprojected
updates is $(I-\alpha_nH_n)(\theta-\theta_n^\circ)$.
The eigenvalues of this symmetric matrix lie in
$[0,1-\alpha_n\mu_n]$, proving contraction by that factor. Iterating
$M_n$ times and applying the triangle inequality to the previous target
gives $e_n\le r_n(e_{n-1}+d_n)$. Unrolling this scalar recurrence
gives (\ref{eq:audit_optimizer_tracking}).

For the bias claim set $d=\theta_n^\circ-\theta_n^*$.
The variational inequality at the constrained minimizer gives
$d^{\mathsf T}\nabla\mathcal J_n(\theta_n^\circ)\le0$.
Since the gradient difference is $H_nd$,
\begin{equation*}
 \mu_n\|d\|^2\le d^{\mathsf T}H_nd
 \le-d^{\mathsf T}\nabla\mathcal J_n(\theta_n^*)
 \le\|d\|\,\|\nabla\mathcal J_n(\theta_n^*)\|.
\end{equation*}
Divide by $\|d\|$ when nonzero; the zero case is immediate.
\end{IEEEproof}
Here $\nabla\mathcal J_n(\theta_n^*)$ measures retained-data
incompatibility in free-coefficient directions. It can be nonzero because
of sensing noise, quadrature, model mismatch, or a physical change within
retained history. A positive empirical curvature caused partly by noise
does not imply a small bias term. At a material jump the target drift can
be large; neither the bound nor the experiment implies instantaneous
tracking. When $H_n$ is singular, the descent result remains valid but
strict contraction and unique identification do not follow.

\subsection{From coefficient error to evaluated derivative error}
For any of these prescribed families, stack its evaluated partials into
one vector and write
\begin{equation}
 \operatorname{vec}F_\theta(z)=a(z)+B(z)\theta.
 \label{eq:audit_partial_map}
\end{equation}
The fixed terms are in $a$; optimization scales are included in $B$.
Equations (\ref{eq:audit_contact_partials}),
(\ref{eq:audit_fourier_partials}), and
(\ref{eq:audit_burgers_partials}) specify this map directly.

\begin{auditproposition}[Evaluated derivative tracking]
\label{prop:audit_partial_error}
Let $W$ be any fixed diagonal weighting of vectorized partial entries and
let $P_U$ be the diagonal coordinate mask, of the same dimension,
that retains only entries depending on unknown coefficients.
At a common context $z$, the same-family coefficient target obeys
\begin{align}
 \|WP_U\operatorname{vec}(F_\theta-F_{\theta^*})\|
 &\le\|WP_UB(z)\|\,\|\theta-\theta^*\|.\label{eq:audit_partial_error}
\end{align}
Under Proposition~\ref{prop:audit_tracking}, this error for $\theta_n$
is bounded by
\begin{equation}
 \|WP_UB(z)\|
 \left(e_n+\frac{\|\nabla\mathcal J_n(\theta_n^*)\|}{\mu_n}\right).
 \label{eq:audit_partial_tracking}
\end{equation}
If the physical differential is not in the family, its discrepancy from
$F_{\theta_n^*}$ must be added to the right-hand side.
\end{auditproposition}
\begin{IEEEproof}
Subtract (\ref{eq:audit_partial_map}) at $\theta$ and $\theta^*$.
The fixed part cancels and the operator-norm inequality gives
(\ref{eq:audit_partial_error}). Use the triangle inequality
$\|\theta_n-\theta_n^*\|\le e_n+
\|\theta_n^\circ-\theta_n^*\|$ and
(\ref{eq:audit_physical_bias}) for the second conclusion. A second
triangle inequality adds any outside-family discrepancy.
\end{IEEEproof}
This calculation explains both the capability and its limitation:
constant learned coefficients can represent rapidly changing evaluated
partials, but successful optimization does not ensure accurate physical
partials without small bias and a suitable family. The norm of $B(z)$
can amplify coefficient error, especially in the Fourier phase column
or contact curvature term. This is why evaluated partials, rather than
coefficient error alone, are reported.
The reported pre-update partials use $\theta_{n-1}$ at the current context.
For those scores, the same bound holds with $e_n$ replaced by
$e_{n-1}+d_n$; this follows before applying the current optimization
steps. Post-update improvement is not substituted for pre-update accuracy.
For contact, the optimizer bound is applied separately to the material
fit and each kinematic row, using their own curvature and step sizes;
the three coefficient-error bounds combine by the Euclidean norm before
applying the full derivative map. A common step size across these fits
is not assumed.

The unknown-only diagnostic excludes the supplied $\rho$, unit blocks,
and fixed zero entries in contact; the $-1$ input partial in Fourier;
and the unit $w_t$ partial in Burgers. With the same physical scaling,
define
\begin{equation}
 e_{F,U}=\frac{\|WP_U\operatorname{vec}(\widehat F-F)\|}
 {\|WP_U\operatorname{vec}F\|}.
 \label{eq:audit_unknown_only}
\end{equation}
This ratio requires a nonzero denominator; zero-denominator cases must
be reported separately. Correct supplied blocks contribute nothing to
the error numerator but increase the full-matrix denominator, so full
relative errors can understate error in learned entries. Both conventions
should be distinguished. Neither pointwise derivative accuracy nor a
small compatibility energy alone establishes closed-loop stability.

\subsection{Information-flow checks and remaining limitations}
The public test file \path{tests/test_information_boundaries.py}
deliberately replaces the clean diagnostic histories and the evaluated
Jacobian functions with corrupted values while preserving the plant,
sensors, references, and random stream. For the three structured learned branches,
the physical trajectories, controls, and learned coefficient sequences
remain identical; the diagnostic outputs change. A positive-control test
changes Fourier's measured rate noise and observes changed learning,
while preserving the first premeasurement command. Another test forbids
access to the true-coefficient configuration when constructing and
updating the Burgers learner. These tests complement source tracing and
the existing finite-difference and current-candidate memory tests. They
do not constitute a proof that every possible future code change or
unexamined branch is free from information leakage.

The present audit found no evaluator-to-learner or evaluator-to-controller
path in the adaptive branches examined. True coefficients generate the
physical plant and simulated measurements, and are used by explicitly
named oracle baselines and evaluators. This is different from supplying
them as learning targets. Nevertheless, the examples assume correct
functional families, physically informed coefficient boxes, selected
known coefficients, and calibrated levels. Fourier additionally supplies
phase and frequency; Burgers supplies exact state/spatial channels and
matched discretization; contact measures force direction. These declared
advantages prevent interpreting the experiments as unrestricted equation
discovery or as validation of every convergence hypothesis of the
continuous-time theory.
